\documentclass[10pt,a4paper]{amsart}
\usepackage[margin=19mm]{geometry}
\usepackage{amsmath,amssymb,mathtools}
\usepackage[scr=rsfs,cal=euler]{mathalfa}
\usepackage{xfrac}
\usepackage[unicode,hidelinks,bookmarks=false]{hyperref}
\newcommand{\ie}{i.e.\ }
\newcommand{\cf}{cf.\ }
\newcommand{\resp}{respectively\ }
\newcommand{\Subsection}{\subsection}
\providecommand{\nobreakdash}{\nobreak\hbox{-}}
\setlength{\emergencystretch}{2em}
\multlinegap0pt
\usepackage{bm}
\usepackage{bbm}
\usepackage{dsfont}
\usepackage{xspace}
\usepackage{cancel}
\usepackage{mathtools}
\usepackage{amsthm}
\makeatletter
\@ifundefined{theorem}{\newtheorem{theorem}{Theorem}}{}
\@ifundefined{lemma}{\newtheorem{lemma}[theorem]{Lemma}}{}
\@ifundefined{proposition}{\newtheorem{proposition}[theorem]{Proposition}}{}
\@ifundefined{remark}{\newtheorem{remark}[theorem]{Remark}}{}
\makeatother
\makeatletter
\def\D {\mathcal{D}}
\def\E{{\mathbb E}}
\def\L {\bar{L}}
\def\N{{\mathbb N}}
\expandafter\ifx\csname proof\endcsname\relax
\newenvironment{proof}[1][Proof.]{\textbf{#1} }{\hfill $\blacksquare$}
\fi
\def\z {\mathbf{z}}
\makeatother
\title[Functional law of large numbers and central limit theorem fo]{Functional law of large numbers and central limit theorem for Crump-Mode-Jagers branching processes}
\author{Ibrahima Dram\'e, Etienne Pardoux}
\date{Source: arXiv:2508.12058v1 (CC BY 4.0)}
\begin{document}
\maketitle

\noindent\emph{Source and licence.} Functional law of large numbers and central limit theorem for Crump-Mode-Jagers branching processes. Version arXiv:2508.12058v1 (CC BY 4.0), distributed under the Creative Commons Attribution 4.0 International licence, as recorded in the arXiv licence metadata for this exact version and shown on the abstract page https://arxiv.org/abs/2508.12058v1. Full rights notice: licenses/arxiv-2508.12058-CC-BY-4.0.txt.\par\smallskip

\noindent\textit{Editorial scope.} Counted unit: the Lemma whose source label is \texttt{MajorG24} in the article (source0.tex lines 574--579, source label \texttt{MajorG24}), delivered together with the article's own complete original proof of it (source0.tex lines 580--598, one \texttt{proof} environment of 19 lines). No mathematics is written, repaired or re-derived: statement and proof are the article's own text line for line. Required context retained uncounted: M3tC (lines 327--332); MomComp (lines 461--471); ecartZ (lines 518--520); ecartU (lines 522--525); ecartEZ (lines 527--532); MajorG (lines 554--560); basic (lines 562--567). Every definition, display and label of those lines is retained unchanged. Omitted from this package and not redistributed: 1--326, 333--460, 472--517, 521--521, 526--526, 533--553, 561--561, 568--573, 599--872. Nothing inside a retained excerpt is omitted or altered: every retained body is delivered exactly as the article's lines are written, so no omission receipt of any kind accompanies this package. Citations: the retained excerpts carry no citation command at all, so no bibliography entry of the article is transcribed and no reference is redistributed. Reference adaptation: none was needed and none was made; every reference inside the delivered unit resolves inside it, so no unresolved reference or citation is produced. Printed slips kept literal: no printed word, symbol, hypothesis or number of the counted statement or of its proof was altered. Layout and class: the article's own class is \texttt{article} with packages including color, amsfonts, amsmath, amssymb, graphicx, mathptmx, subfigure, float, parskip, hyperref, dsfont; the wrapper is the lane's pinned \texttt{amsart} editorial wrapper at 10pt on A4, which restates the theorem-like environments the retained text needs and adds the article's own macro definitions that the retained text needs (\texttt{\textbackslash D}, \texttt{\textbackslash E}, \texttt{\textbackslash L}, \texttt{\textbackslash N}, \texttt{\textbackslash z}), with extra wrapper packages (bm, bbm, dsfont, xspace, cancel, mathtools). The delivered numbering is the unit's own. Assets: no retained span contains a figure environment, a graphics-inclusion command, a PDF-page inclusion, a table or a TikZ picture; no image file is copied into this package, the package declares no asset dependency and no image file is redistributed with it. Licence: version arXiv:2508.12058v1 of this article is distributed under the Creative Commons Attribution 4.0 International licence, as recorded in the arXiv licence metadata for this exact version and shown on the abstract page https://arxiv.org/abs/2508.12058v1. A full rights notice is delivered at licenses/arxiv-2508.12058-CC-BY-4.0.txt. Characters: every retained excerpt is pure ASCII, so the delivered engine, XeLaTeX with its default Latin Modern text font, reports no missing character.
 \begin{proposition}\label{M3tC}
Suppose that Assumption ${(\bf H_2)}$ is satisfied. For any $T>0$, there exists $C$ such that 
 \begin{align*}
\sup_{0<t<T} \E([Z(t)]^2)\le C.
\end{align*} 
  \end{proposition} 

\begin{proposition}\label{MomComp}
Let $(E, \mathcal{E})$ be a measurable space. Let $\mathbf{Q}$ be a Poisson random measure on $(E, \mathcal{E})$ with mean $\mu$, and $\bar{\mathbf{Q}}$ its associated 
compensated measure. If $f_i\in L^1(\mu) \cap L^4(\mu)$, 
\begin{align*}
\E[(\bar{\mathbf{Q}}(f))^2]=\mu(f^2) \quad \mbox{and} \quad \E[(\bar{\mathbf{Q}}(f))^4]&=\mu(f^4) + 3[\mu(f^2)]^2, 
\end{align*}
and if $f_i\in L^1(\mu) \cap L^4(\mu)$, $i=1,2$, 
 \begin{align*}
\E \left[(\bar{\mathbf{Q}}(f_1))^2\bar{\mathbf{Q}}(f_2)^2 \right]=\mu(f_1^2f_2^2)+\mu(f_1^2)\mu(f_2^2)+2[\mu(f_1f_2)]^2\,.
\end{align*}
\end{proposition}

 \begin{align}\label{ecartZ}
 Z(t)-Z(s)= -{\bf1}_{s<\eta\le t}+ \int_s^t \L\, M_1(t-r)b(r)dr+\int_0^s \L\,[M_1(t-r)-M_1(s-r)]b(r)dr + U(t)-U(s),
  \end{align}

 \begin{align}\label{ecartU}
U(t)-U(s)=\int_s^t\int_0^{b(r)}\int_\D \sum_{i=1}^\ell z_i(t-r)\bar{\mathcal{M}}(dr,du,d\z)
 +\int_0^s\int_0^{b(r)}\int_\D \sum_{i=1}^\ell [z_i(t-r)-z_i(s-r)]\bar{\mathcal{M}}(dr,du,d\z) \, .
 \end{align}

  \begin{lemma}\label{ecartEZ}
For $0< s<t$, we have
 \begin{align*}
 |\E[Z(t)-Z(s)]|= |M_1(t)-M_1(s)|\le G(t)-G(s).
 \end{align*}
 \end{lemma}

\begin{remark}\label{MajorG}
For $0<s<t$, it is easy to check that (recall that $s<t\le T$ and $T$ is fixed)
 \begin{align*}
(1)& \quad \int_0^s[G(t-r)-G(s-r)]dr \le \int_0^t G(r) dr-  \int_0^sG(r)dr  \le C(t-s)\\
(2) &\quad  G(t)-G(s)  \le C(t-s)^\alpha.
 \end{align*}
 \end{remark}

  \begin{lemma}\label{basic}
  Let $n\in\N$, $\xi$ be an $\N$--valued r.v. and $\{X_i, i\ge1\}$ a sequence of identically distributed rv's
  having moment of order $n$,
  which is globally independent of $\xi$. Then
  \[ \E\left[\left(\sum_{i=1}^\xi X_i\right)^n\right]\le \E[\xi^n]\E[X_1^n]\,.\]
  \end{lemma}

  \begin{lemma}\label{MajorG24}
For $0< s<t$, we have
 \begin{align*}
 \E[|Z(t)-Z(s)|^2]\le G(t)-G(s)\quad \mbox{and} \quad \E[|Z(t)-Z(s)|^4]\le G(t)-G(s).
 \end{align*}
 \end{lemma}

\begin{proof} From \eqref{ecartZ}, \eqref{ecartU} and Lemma \ref{basic}, we have
  \begin{align*}
  \E[|Z(t)-Z(s)|^2]&\le 5[F(t)-F(s)]+5\L^2(t-s)\int_s^t\E[b(r)^2]dr  +5\int_s^t \E[L^2]\E[(Z(t-r))^2] \bar{b}(r)dr  \\
  &\quad + 5s\L^2\int_0^s(M_1(t-r)-M_1(s-r))^2\E[b^2(r)]dr +5\int_0^s \E[L^2]\E[|Z(t-r)-Z(s-r)|^2]\bar{b}(r)dr.
  \end{align*}
However, using Lemma \ref{ecartEZ} and $(1)$ of Remark \ref{MajorG}, we note that 
  \begin{align*}
  5s\L^2\int_0^s(M_1(t-r)-M_1(s-r))^2\E[b^2(r)]dr \le C\int_0^s|M_1(t-r)-M_1(s-r)|dr \le C(t-s).
  \end{align*}
  Combining the last two estimates, our assumptions, Proposition \ref{M3tC} and the definition of $G$, we deduce that 
  \begin{align*}
  \E[|Z(t)-Z(s)|^2]&\le  G(t)-G(s)+C\int_0^s \E[|Z(t-r)-Z(s-r)|^2]dr,
  \end{align*}
  Hence the first assertion follows from Gronwall's Lemma applied to $\varphi(u)=\E[|Z(t-s+u)-Z(u)|^2)$, $0\le u\le s$.
% Using the same argument as in Lemma \ref{ecartEZ}, we deduce that
%  \[ \E[|Z(t)-Z(s)|^2]\le G(t)-G(s),\]
%  provided we have redefined $G$ by multiplying it by an appropriate constant. `
  The second assertion is easily deduced by a similar argument, exploiting the fourth moment estimate in Proposition \ref{MomComp}. 
 \end{proof}

\end{document}
